% GENERATED FILE. Edit canonical lessons, then run npm run generate:books.
% canonical-source-hash: fnv1a64:8c5d665b8f8fa2a2
% canonical-lessons: FR-C384-expirer, FR-C384-valable, FR-C384-valide, FR-C384-renouveler, FR-C384-renouvellement

\chapter{Des papiers en règle}
\label{ch:fr-des-papiers-en-regle}
\hlchaptermodality{384}

\begin{chapteropening}
\textbf{By the end of this chapter:} \emph{I can keep my papers in order: say when something expires, whether it is still valid, and renew it.}

Complete the last lesson of chapter 384: I can keep my papers in order: say when something expires, whether it is still valid, and renew it.
\end{chapteropening}

\section[expirer]{expirer — to expire; to breathe out}

\label{lesson:FR-C384-expirer}



\begin{quote}
Before the new one: say the French for the admission, then the French for to exclude.
\end{quote}

\subsection*{You'll want to know: expirer}
\textbf{expirer} — \textquotedblleft{}to expire; to breathe out\textquotedblright{}.

\begin{quote}
Mon passeport expire au mois de mai.
\end{quote}

\begin{quote}
\emph{My passport expires in May.}
\end{quote}

Of breathing: \emph{inspirer}, then \emph{expirer}.

\begin{cousinweb}
\emph{expirer} is a learned borrowing from Latin \emph{exspīrāre} \textquotedblleft{}to breathe out, to die away\textquotedblright{} (EtymDB; Lewis \& Short), built on \emph{spīrāre} \textquotedblleft{}to breathe\textquotedblright{}.
\end{cousinweb}

\begin{grammarlens}[title={What you've built}]
One more word for telling, explaining and describing at B1.
\end{grammarlens}

\subsection*{Your turn}
\begin{itemize}
\raggedright
  \item \emph{Say it:} \emph{expirer}
  \item \emph{Say it:} \emph{expirer}, once more
  \item \emph{Say it:} a sentence of your own with \emph{expirer}
  \item \emph{Recall:} say the French for the admission, then the French for to exclude, then say \emph{expirer} again
\end{itemize}

\begin{tcolorbox}[breakable,colback=teal!4,colframe=teal!35!black,title={Before you move on}]
What does expirer mean? (\textquotedblleft{}To expire; to breathe out\textquotedblright{}.) Say it once more.
\end{tcolorbox}
\section[valable]{valable — acceptable, valid; worthwhile}

\label{lesson:FR-C384-valable}



\begin{quote}
Before the new one: say the French for to exclude, then the French for to expire.
\end{quote}

\subsection*{You'll want to know: valable}
\textbf{valable} — \textquotedblleft{}acceptable, valid; worthwhile\textquotedblright{}.

\begin{quote}
Ce billet est valable une semaine.
\end{quote}

\begin{quote}
\emph{This ticket is valid for a week.}
\end{quote}

\emph{Une excuse valable} — a good excuse. One form for both genders.

\begin{cousinweb}
\emph{valable} is built from \emph{valoir} \textquotedblleft{}to be worth\textquotedblright{} and \emph{-able} (EtymDB). \emph{valoir} is inherited from Latin \emph{valēre} \textquotedblleft{}to be strong, to be worth\textquotedblright{} (Lewis \& Short).
\end{cousinweb}

\begin{grammarlens}[title={What you've built}]
One more word for telling, explaining and describing at B1.
\end{grammarlens}

\subsection*{Your turn}
\begin{itemize}
\raggedright
  \item \emph{Say it:} \emph{valable}
  \item \emph{Say it:} \emph{valable}, once more
  \item \emph{Say it:} a sentence of your own with \emph{valable}
  \item \emph{Recall:} say the French for to exclude, then the French for to expire, then say \emph{valable} again
\end{itemize}

\begin{tcolorbox}[breakable,colback=teal!4,colframe=teal!35!black,title={Before you move on}]
What does valable mean? (\textquotedblleft{}Acceptable, valid; worthwhile\textquotedblright{}.) Say it once more.
\end{tcolorbox}
\section[valide]{valide — valid; fit, able-bodied}

\label{lesson:FR-C384-valide}



\begin{quote}
Before the new one: say the French for to expire, then the French for acceptable.
\end{quote}

\subsection*{You'll want to know: valide}
\textbf{valide} — \textquotedblleft{}valid; fit, able-bodied\textquotedblright{}.

\begin{quote}
Il faut une pièce d'identité valide.
\end{quote}

\begin{quote}
\emph{A valid identity document is needed.}
\end{quote}

Of a person: fit, able-bodied.

\begin{cousinweb}
\emph{valide} is a learned borrowing from Latin \emph{validus} \textquotedblleft{}strong, powerful\textquotedblright{} (EtymDB; Lewis \& Short), from \emph{valēre} \textquotedblleft{}to be strong\textquotedblright{}.
\end{cousinweb}

\begin{grammarlens}[title={What you've built}]
One more word for telling, explaining and describing at B1.
\end{grammarlens}

\subsection*{Your turn}
\begin{itemize}
\raggedright
  \item \emph{Say it:} \emph{valide}
  \item \emph{Say it:} \emph{valide}, once more
  \item \emph{Say it:} a sentence of your own with \emph{valide}
  \item \emph{Recall:} say the French for to expire, then the French for acceptable, then say \emph{valide} again
\end{itemize}

\begin{tcolorbox}[breakable,colback=teal!4,colframe=teal!35!black,title={Before you move on}]
What does valide mean? (\textquotedblleft{}Valid; fit, able-bodied\textquotedblright{}.) Say it once more.
\end{tcolorbox}
\section[renouveler]{renouveler — to renew}

\label{lesson:FR-C384-renouveler}



\begin{quote}
Before the new one: say the French for acceptable, then the French for valid.
\end{quote}

\subsection*{You'll want to know: renouveler}
\textbf{renouveler} — \textquotedblleft{}to renew\textquotedblright{}.

\begin{quote}
Je dois renouveler ma carte d'identité.
\end{quote}

\begin{quote}
\emph{I have to renew my identity card.}
\end{quote}

\emph{Je renouvelle}, with \emph{ll} before a silent ending.

\begin{cousinweb}
\emph{renouveler} is built from \emph{re-} and Old French \emph{noveler} \textquotedblleft{}to make new\textquotedblright{}, from Latin \emph{novellus} \textquotedblleft{}young, new\textquotedblright{}, a small form of \emph{novus} \textquotedblleft{}new\textquotedblright{} (Lewis \& Short).
\end{cousinweb}

\begin{grammarlens}[title={What you've built}]
One more word for telling, explaining and describing at B1.
\end{grammarlens}

\subsection*{Your turn}
\begin{itemize}
\raggedright
  \item \emph{Say it:} \emph{renouveler}
  \item \emph{Say it:} \emph{renouveler}, once more
  \item \emph{Say it:} a sentence of your own with \emph{renouveler}
  \item \emph{Recall:} say the French for acceptable, then the French for valid, then say \emph{renouveler} again
\end{itemize}

\begin{tcolorbox}[breakable,colback=teal!4,colframe=teal!35!black,title={Before you move on}]
What does renouveler mean? (\textquotedblleft{}To renew\textquotedblright{}.) Say it once more.
\end{tcolorbox}
\section[le renouvellement]{le renouvellement — the renewal}

\label{lesson:FR-C384-renouvellement}



\begin{quote}
Before the new one: say the French for valid, then the French for to renew.
\end{quote}

\subsection*{You'll want to know: le renouvellement}
\textbf{le renouvellement} — \textquotedblleft{}the renewal\textquotedblright{}.

\begin{quote}
Le renouvellement du passeport prend un mois.
\end{quote}

\begin{quote}
\emph{Renewing a passport takes a month.}
\end{quote}

Spelled with \emph{ll}.

\begin{cousinweb}
\emph{renouvellement} is built from \emph{renouveler} and \emph{-ment} (EtymDB). The verb rests on Latin \emph{novellus} \textquotedblleft{}new\textquotedblright{}, from \emph{novus} \textquotedblleft{}new\textquotedblright{} (Lewis \& Short).
\end{cousinweb}

\begin{grammarlens}[title={What you've built}]
One more word for telling, explaining and describing at B1.
\end{grammarlens}

\subsection*{Your turn}
\begin{itemize}
\raggedright
  \item \emph{Say it:} \emph{le renouvellement}
  \item \emph{Say it:} \emph{le renouvellement}, once more
  \item \emph{Say it:} a sentence of your own with \emph{le renouvellement}
  \item \emph{Recall:} say the French for valid, then the French for to renew, then say \emph{le renouvellement} again
\end{itemize}

\begin{tcolorbox}[breakable,colback=teal!4,colframe=teal!35!black,title={Before you move on}]
What does le renouvellement mean? (\textquotedblleft{}The renewal\textquotedblright{}.) Say it once more.
\end{tcolorbox}
